\begin{frame}
	\frametitle{Pseudo-tested Methods}
	\begin{center}
    \huge{Pseudo-tested Methods}
  \end{center}
\end{frame}

\subsection{Defintion}
\begin{frame}
  \frametitle{Definition}
    \begin{center}
      \textbf{What is a Pseudo-tested Method?}

      \vspace{10mm}
      \includegraphics{images/passing}

      \vspace{10mm}
      \textbf{\textit{Def}}: It will never fail.
    \end{center}
\end{frame}

\subsection{Detection}
\begin{frame}
  \frametitle{Detection}
    \begin{center}
      \textbf{How Can We Detect Pseudo-tested Methods}

			\vspace{20mm}
			\textbf{It is harder than you think!}

    \end{center}
\end{frame}

\begin{frame}[fragile]
\frametitle{Example of a Pseudo-tested method}
	\begin{figure}[t!, scale = .75]
	\begin{lstlisting}[language = Python, basicstyle=\small, backgroundcolor = \color{lightgray}]
numbers.py:
def numberOrder(n):
  numbersSorted = sorted(n)
  return numbersSorted


test_numbers.py:
def test_numbers_ordered():
  numbers = set([2,4,3,1])
  sortedNumbers = set([1,2,3,4])
  orderedNumbers = numberOrder(numbers)
  assert numbers == sortedNumbers
	\end{lstlisting}
	\end{figure}
\end{frame}
